\subsection{indexation iconographique}

\begin{itemize}

\item Définition
«L'indexation iconographique consiste à décrire le contenu des images à l'aide de descripteurs afin de faciliter l'accès aux œuvres, sans cependant dispenser sa consultation.»
\footcite[p.~3]{husson2025}

\item Outils 

"When it comes to processing image-based records in archives, supervised learning models can be	helpful for identification and classification, while unsupervised learning	models are best suited for processing large amounts of data fo creating	generic descriptions, and highlighting possible connections between records."
\footcite{fewster}

\item cas d'usage 

Computer–Aided Metadata generation for Photo Archives Initiative (CAMPI) \footcite{fewster} production de tags sur des images et recommandation d'images similaires 

\subitem détection d'objets 

YOLO (you only look once)
Algorithme qui utilise CNNs pour entourer des objets dans des images et les classer dans des catégories 
il s'agit d'un single shot object detector : il se base sur une seule analyse de l'image dans sa globalité 
facial recognition, scene analysis, object tracking \footcite{fewster}

two shot objects detector : multiple image reviews 
Exemple : Mask	Region-based CNN (R-CNN) algorithm. "uses deep neural networks to	identify regions of interest in the photograph and build boxes around objects	found throughout the image" \footcite{fewster}
Cette manière de faire est appelée "semantic segmentation" : "identifies semantic objects	without specifying how many instances of that object are present in the	image" \footcite{fewster}
Par exemple, s'il y a 5 personnes dans une image, la segmentation sémantique les identifie comme un même objet sémantique dans l'image.

	Instance segmentation : "identifies all the semantic objects in	an image and segments them into specific instances to accurately distinguish	between each object of a similar semantic category" (c'est ce que fait Mask RCNN)\footcite{fewster}


\end{itemize}